\section*{\centering INTRODUCTION GÉNÉRALE}
\addcontentsline{toc}{section}{INTRODUCTION GÉNÉRALE}

\indent Les robots de service autonomes occupent une place croissante dans les environnements professionnels accueil, guidage, télépresence grâce à la maturité des technologies de navigation autonome et d'interaction homme-robot. Ces robots sont cependant généralement livrés avec un écosystème logiciel fermé : application propriétaire, protocole non documenté, absence d'API publique. Cette fermeture limite fortement la capacité des utilisateurs à adapter leur comportement à des besoins spécifiques.

C'est dans ce contexte que s'inscrit le robot de réception mobile \textbf{CIOT TY1251D-03195}, mis à disposition par HESTIM Engineering \& Business School dans le cadre d'un stage de fin d'année. Composé d'un châssis de navigation sous ROS et d'une tête Android assurant l'interface visiteur, ce robot dépend entièrement d'une application propriétaire pour laquelle aucune documentation n'est disponible. Le projet \textbf{CYBEL}, objet de ce rapport, vise à concevoir une plateforme logicielle indépendante permettant de superviser, piloter et faire interagir ce robot sans recourir à l'application d'origine.

La problématique du stage est donc double : comment établir une communication fiable avec un système robotique fermé dont ni le protocole, ni les points d'accès réseau ne sont connus, et comment construire à partir de cette compréhension reconstruite une solution de contrôle robuste et extensible ? Le stage, d'une durée de trois à quatre mois, poursuit ainsi deux objectifs : procéder à une rétro-ingénierie méthodique du protocole de communication interne du robot, puis développer sur cette base une plateforme de commande complète, incluant une interface opérateur, une interface visiteur et des fonctionnalités d'interaction tactile et vocale. La démarche adoptée repose sur une approche itérative et prudente, où chaque hypothèse sur le fonctionnement du robot est systématiquement vérifiée avant d'être considérée comme acquise.

Le présent rapport restitue cette démarche en sept chapitres : présentation de l'entreprise d'accueil, analyse du besoin et étude préalable, état de l'art, analyse et conception du système CYBEL, méthodologie et planification, réalisation et implémentation, puis validation, résultats et analyse critique, avant une conclusion générale dressant le bilan du stage et ses perspectives.
